\section{Selmer coranks and fixed local data}\label{sec:preliminaries}

For an elliptic curve $A/\Q$, put
\[
 r(A)=\rank_{\Z}A(\Q),\qquad
 a(A)=\ord_{s=1}L(A,s),\qquad
 s_2(A)=\corank_{\Z_2}\Sel_{2^\infty}(A/\Q).
\]
The Selmer group uses the local Kummer image at every place.
The analytic continuation and functional equation are supplied by modularity
\cite{BCDT2001}. Write $\varepsilon(A)\in\{1,-1\}$ for the root number.

\begin{lemma}[The three ranks]\label{pr:ranks}
There is an exact sequence
\[
0\longrightarrow A(\Q)\otimes\Q_2/\Z_2
\longrightarrow \Sel_{2^\infty}(A/\Q)
\longrightarrow \Sha(A/\Q)[2^\infty]\longrightarrow0.
\]
Consequently
\[
s_2(A)=r(A)+\corank_{\Z_2}\Sha(A/\Q)[2^\infty].
\]
If $a(A)\in\{0,1\}$, then $r(A)=a(A)$ and the whole group
$\Sha(A/\Q)$ is finite. In that case $s_2(A)=a(A)$.
\end{lemma}

\begin{proof}
The Kummer sequences for multiplication by $2^n$, with the local
Kummer conditions imposed, give
\[
0\longrightarrow A(\Q)/2^n A(\Q)
\longrightarrow \Sel_{2^n}(A/\Q)
\longrightarrow \Sha(A/\Q)[2^n]\longrightarrow0.
\]
Their direct limit is the displayed sequence. Taking $\Z_2$-coranks proves
the identity. The analytic-rank-zero and analytic-rank-one assertions are
the Gross--Zagier--Kolyvagin theorem
\cite{GrossZagier1986,Kolyvagin1990}; see also the precise unrestricted
forward statement in \cite[\S1]{Skinner2020}.
\end{proof}

No finite $2$-Selmer dimension will be substituted for $s_2$.
In particular, $s_2(A)=1$ does not supply a rational point of infinite order:
the exact sequence allows a contribution from the divisible part of $\Sha$.

\begin{lemma}[Parity and change of curve]\label{pr:invariance}
For every elliptic curve $A/\Q$,
\[
(-1)^{s_2(A)}=\varepsilon(A).
\]
A rational isogeny preserves $a$, $r$, and $s_2$. If
$K=\Q(\sqrt{k})$ is quadratic, then
\[
s_2(A/K)=s_2(A/\Q)+s_2(A^{(k)}/\Q).
\]
\end{lemma}

\begin{proof}
The $2$-parity assertion is due to Monsky \cite{Monsky1996};
see also \cite[Theorems~1.4 and~4.19]{DokchitserDokchitser2010},
which states the Selmer-parity result for every prime and credits
Monsky for the case $2$.
For an isogeny and its dual, both composites are multiplication by the
degree. The corresponding maps on rational points and on Selmer groups
respect the Kummer conditions. After rationalizing their Tate modules,
the maps are inverse up to a nonzero scalar; hence the ranks and coranks
agree. The equality of $L$-functions is isogeny invariance.

Over $K$, identify $A^{(k)}$ with $A$, and let $c$ be the nontrivial
element of $\Gal(K/\Q)$. Restriction and the two conjugation
projections $1+c$ and $1-c$ define maps between the sum of the two
Selmer groups over $\Q$ and the Selmer group over $K$ whose composites are
multiplication by $2$. Functoriality of local Kummer images gives these
maps at every place. Rationalization therefore gives the stated corank
identity. This argument does not assert an integral direct-sum
decomposition of the Selmer groups.
\end{proof}

Fix a curve $E/\Q$ of conductor $N$. A finite support $S$ always contains
the primes dividing $2N$. It may be enlarged by finitely many primes when
the fixed construction is prepared. A \emph{local filter} consists of a
real sign and unit squareclasses at $S$ for the variable fundamental
discriminant. The filter is imposed on every relevant Fourier index,
not just the squarefree indices. Its root number is constant.
All variable primes lie outside $S$.

For an odd prime $p$, set
\[
p^*=\left(\frac{-1}{p}\right)p.
\]
Nonempty products of distinct $p^*$ are odd fundamental discriminants.
We also admit the empty product $h=1$, which denotes the untwisted
curve rather than a quadratic field. For a squarefree product $h$ of
such factors, set
\[
 t_p=\dim_{\F_2} E[2]^{\operatorname{Fr}_p},
 \qquad
 w(h)=\frac12\sum_{p\mid h}t_p.
\]
In the rational-two-torsion range, $t_p$ is $1$ or $2$.
Whenever the $2$-primary group below is finite, write
\[
s(h)=\length_{\Z_2}\Sha(E^{(h)}/\Q)[2^\infty].
\]
Thus $s(h)$ is a finite length, whereas $s_2(E^{(h)})$ is a corank.
In particular $s(h)$ is defined when $s_2(E^{(h)})=0$: the exact
sequence in Lemma~\ref{pr:ranks} and finite generation of the dual
Selmer group make its $2$-primary quotient finite. This conclusion
concerns the $2$-primary group; the low-analytic-rank theorem supplies
finiteness of the whole group after analytic nonvanishing is proved.
Let $v$ be a valuation above $2$, normalized by $v(2)=1$.
After the coefficient field has been fixed, its value group is discrete.

The notation $L_*(h)$ denotes $L(E^{(h)},1)\sqrt{|h|}$ divided by the
fixed real or imaginary period appropriate to the filter.
The notation $\mathcal L(h)$ in Section~\ref{cy:section} denotes
this same normalized value.
Every modular parametrization used to form Heegner points is normalized
to send the cusp $\infty$ to the identity of the elliptic curve.
Changing among fixed nonzero modular parametrizations, isogenous curves,
or period conventions contributes only bounded constants to valuations.
Every such dependence is fixed before a binary cube is allowed to grow.
Dependence on a growing number of prime factors will always be stated
explicitly.

The following ledger records the order of choices used in the proof.
It distinguishes bounds from the sizes and precisions needed to realize
the auxiliary primes.
\begin{center}
\small
\begin{tabular}{@{}>{\raggedright\arraybackslash}p{0.29\textwidth}>{\raggedright\arraybackslash}p{0.65\textwidth}@{}}
\toprule
Data or bound & Allowed dependence and order of choice \\
\midrule
Fixed arithmetic constants & The curve, its parametrization and periods,
the fixed part of the support, and the chosen local classes; these are fixed before
the varying prime support or cube dimension. \\
Varying quadratic companions & The same data and a bound $A$ for
the number of prime factors of the companion discriminant; its prime
sizes do not enter. \\
Fixed-companion transfer & The companion field is fixed before the
cube dimension. Constants may depend on its fixed base Selmer data,
as specified in Section~\ref{rd:section}. \\
Finite realization & Choose the cube dimension and finite networks,
then the required coefficient precision and a sufficiently late
arithmetic stage. Prime sizes and stage thresholds may depend on
these choices. \\
\bottomrule
\end{tabular}
\end{center}
